\section{Introduction}
\label{sec:introduction}

\subsection{Background, Motivation and Literature Review}
\label{sec:background_and_motivation}

Reaching a net zero economy requires a massive undertaking in electrification of the energy system. The majority of the effort occurs on three fronts:
\begin{enumerate}
    \item shift towards renewable power generation, e.g., wind and photovoltaic generation, 
    \item electrification of energy demand, e.g., transport, heating and industrial processes, and
    \item expansion of the transmission system to facilitate the increase in electric load and generation, and the geographic shift of latter.
\end{enumerate}

In the context of power quality this brings about two major shifts. First, renewable energy sources and new types of electrical loads are almost exclusively power electronic interfaced assets. Power electronic interfaced assets behave fundamentally different compared to classical electrical machines, or resistive and inductive demand. One of the major concerns is the increase in injected current at non-fundamental frequencies, i.e., harmonics. Second, due to societal pressure, in high density areas, underground cables have become the preferred option for transmission system expansion in favor of overhead lines. As a consequence, the impedance of the transmission system shifts from a purely inductive impedance, i.e., linearly increasing with frequency, to an impedance containing resonances, i.e., going to zero or infinity at certain frequencies. Sun~et~al. highlight that the traditional harmonic studies performed at the initial planning stage may no longer be appropriate in this context, ~\cite{sun_distribution_2022}. Considering the aforementioned transitions, this paper aims at answering the question: “what is the maximum harmonic distortion, e.g., number of power electronic converters, that can be reliably hosted in a given power system, with increasing degrees of underground cable connections?”. In literature, this concept is referred to as the harmonic hosting capacity of a power system~\cite{bollen_hosting_2017,koirala_hosting_2022}. 

There are several standards and, more recently, limited academic literature presenting methods to determine the harmonic hosting capacity of a power system, and differing in a number of features, including network representation. The methods in the relevant standards, e.g., IEEE519-2022~\cite{noauthor_ieee519_2022} and~IEC61000-3-6:2008~\cite{noauthor_iec61000-3-6_2008}, use a nodal representation of the network, i.e., calculate the harmonic hosting capacity for each individual bus, (partially) omitting the interaction between different busses. The standard IEC61000-3-6:2008, at pg.~35, partially compensates for omitting this interaction by the introduction of influence factors, adding significant complexity. Furthermore, in the presence of series resonances, i.e., where the nodal impedance goes to zero, the method described in the standard is unable to ensure compliance with the prescribed individual harmonic voltage limits. The standard IEC61000-3-6:2008, at pg.~23, partially compensates the series resonance effects by disregarding the actual harmonic impedance for around the series resonance frequency, and replacing it with the fundamental impedance scaled with the harmonic number. 

Recent literature shifted towards a full network representation in the context of harmonic hosting capacity. Sun et al. present a non-linear optimization method, including harmonic voltage constraints, to determine the hosting capacity for distributed generation in distribution networks~\cite{sun_distribution_2012},~\cite{sun_incorporating_2012}. Sakar et al. present a monte-carlo power flow method to determine the harmonic hosting capacity for large-scale photovoltaic plants~\cite{sakar_integration_2018} whereas Carmelito presents a similar approach to determine the harmonic hosting capacity for electric vehicles~\cite{carmelito_hosting_2023}. In constrast to Sakar et al., Carmelito considers an unbalanced power flow rather than a balanced one.

\subsection{Contribution and Outline}
\label{sec:contribution_and_outline}

In this paper, a holistic one-shot optimization method is presented to exactly and tractably determine the lower bound of the steady-state harmonic hosting capacity for a phase-balanced three-phase power system. The problem is formulated from the perspective of the power system operator, i.e., without knowledge of the (future) assets of its costumers. Similar to the methods described in the relevant standards, the harmonic hosting capacity is determined in terms of harmonic current magnitude. The models and case study presented in the paper are made available as part of the \textsc{Julia}-package:~\textsc{HarmonicPowerModels.jl}, at:~\href{https://github.com/timmyfaraday/HarmonicPowerModels.jl}{https://github.com/timmyfaraday/HarmonicPowerModels.jl}.

The remainder of the paper is organized as follows: Section~\ref{sec:mathematical_framework} presents the mathematical framework of the methods in the relevant standards and the novel optimization method. Section~\ref{sec:case_study} shows results for the IEEE-30 bus system. Section~\ref{sec:conclusion} concludes the paper.